\documentclass[11pt]{article}
\usepackage[a4paper,margin=25mm]{geometry}
\usepackage{amsmath,amssymb,amsthm}
\usepackage[hidelinks]{hyperref}
\newtheorem{theorem}{Theorem}
\newtheorem{proposition}[theorem]{Proposition}
\theoremstyle{remark}
\newtheorem*{remark}{Remark}
\title{Cycle Graphs Exceed the Conjectured Cycle-Count Bound:\\
A Disproof of Conjecture 00000004413}
\author{GitHub: gaochengzhecpu}
\date{}
\begin{document}
\maketitle
\begin{abstract}
The source asserts that in a graph limit with degree bound $d$ the
measure of cycles of length $\ell$ is at most $(d-1)^\ell/(2\ell)$, and
that this bound is attained and optimal for the $d$-regular tree. For
$d=2$ the bound is false: the cycle graph $C_\ell$ has degree $2$ and
cycle measure $1/\ell>1/(2\ell)$, for every $\ell\ge3$. This is
verified in Lean for Mathlib's graphs. For $d\ge3$ the bound holds but
is neither attained nor optimal, and for no $d\ge2$ does the tree attain
it; these two facts are proved here on paper. An earlier submission
with the triangle as example was withdrawn by its author; it is
discussed below.
\end{abstract}

\section*{Reading of the source}
A bounded-degree graph limit (Benjamini--Schramm limit, or graphing)
is a random rooted graph $(G,o)$ with all degrees at most $d$ that
satisfies the mass-transport principle. Every finite graph $G$ with a
uniformly random root is such an object; it is the limit of the
constant sequence $G,G,\dots$. For a finite graph the \emph{measure of
cycles of length $\ell$} is
\[
 c_\ell(G)=\frac{\#\{\text{cycles of length }\ell\text{ in }G\}}{|V(G)|},
\]
and for a general limit it is $\mathbb E[N_\ell(G,o)]/(2\ell)$, where
$N_\ell(G,o)$ is the number of sequences
$(o=v_0,v_1,\dots,v_{\ell-1})$ of distinct vertices with $v_i$ adjacent
to $v_{i+1}$ for all $i$ modulo $\ell$. The two agree for finite
graphs, because each cycle of length $\ell\ge3$ gives exactly $2\ell$
such sequences in total, one for each starting vertex and direction.
A cycle is a cycle in the graph-theoretic sense: its vertices are
distinct.

The source makes two claims for a degree bound $d$ and a length
$\ell$:
\begin{enumerate}
\item[(i)] $c_\ell\le(d-1)^\ell/(2\ell)$ for every graph limit with
degree bound $d$;
\item[(ii)] this bound is attained, and is optimal, for the
$d$-regular tree.
\end{enumerate}
Other normalisations of the measure are larger: the expected number of
$\ell$-cycles through the root is $\ell\,c_\ell$. So a violation of (i)
for $c_\ell$ is a violation for them as well.

\section*{The disproof}
\begin{theorem}\label{thm:cycle}
Let $\ell\ge3$. The cycle graph $C_\ell$ has all degrees equal to $2$
and $c_\ell(C_\ell)=1/\ell$. Hence
\[
 c_\ell(C_\ell)=\frac1\ell>\frac1{2\ell}=\frac{(2-1)^\ell}{2\ell},
\]
and claim (i) is false for $d=2$ and every $\ell\ge3$.
\end{theorem}
\begin{proof}
Take the vertex set $\mathbb Z/\ell$ with $i$ adjacent to $i\pm1$.
Since $\ell\ge3$, the two neighbours $i-1$ and $i+1$ of $i$ are
different, so every degree is $2$. The graph $C_\ell$ is itself a cycle
of length $\ell$, and it is the only one: a cycle of length $\ell$ has
$\ell$ edges and $C_\ell$ has only $\ell$ edges. So
$c_\ell(C_\ell)=1/\ell$.

For the formal proof only the inequality $c_\ell(C_\ell)\ge1/\ell$ is
needed. The $\ell$ rotations $i\mapsto i+k$ and the $\ell$ reflections
$i\mapsto k-i$ are injective homomorphisms $C_\ell\to C_\ell$. They are
$2\ell$ different maps: a rotation is determined by the image $k$ of
$0$, a reflection likewise, and if $i+k=k'-i$ for all $i$, then $k=k'$
and $2=0$ in $\mathbb Z/\ell$, which is false for $\ell\ge3$. So there
are at least $2\ell$ rooted oriented cycles and
$c_\ell(C_\ell)\ge 2\ell/(2\ell\cdot\ell)=1/\ell$.
\end{proof}

\begin{proposition}\label{prop:upper}
Let $(G,o)$ be a rooted graph with all degrees at most $d$, where
$d\ge1$, and let $\ell\ge3$. Then $N_\ell(G,o)\le d(d-1)^{\ell-2}$.
Hence every graph limit with degree bound $d$ satisfies
\[
 c_\ell\le\frac{d(d-1)^{\ell-2}}{2\ell}.
\]
For $d\ge3$ this is at most $\frac34\cdot\frac{(d-1)^\ell}{2\ell}$.
\end{proposition}
\begin{proof}
Build a sequence $(o=v_0,\dots,v_{\ell-1})$ step by step. There are at
most $d$ choices for $v_1$. For $2\le i\le\ell-1$ the vertex $v_i$ is a
neighbour of $v_{i-1}$ different from $v_{i-2}$, so there are at most
$d-1$ choices. This gives at most $d(d-1)^{\ell-2}$ sequences. Taking
expectations and dividing by $2\ell$ gives the bound for $c_\ell$.
Finally $d(d-1)^{\ell-2}=\frac{d}{(d-1)^2}(d-1)^\ell$, and
$d/(d-1)^2\le3/4$ for $d\ge3$.
\end{proof}

\paragraph{Claim (ii) fails for every $d\ge2$.}
A tree has no cycle, so the $d$-regular tree has $c_\ell=0$, while
$(d-1)^\ell/(2\ell)>0$ for $d\ge2$. So the tree does not attain the
bound. The bound is not optimal either: for $d=2$ it is false by
Theorem~\ref{thm:cycle}, and for $d\ge3$ no graph limit comes within
the factor $3/4$ of it, by Proposition~\ref{prop:upper}. For $d=2$ the
correct optimal bound is $1/\ell$, attained by $C_\ell$, again by
Proposition~\ref{prop:upper}.

\begin{remark}
The expression $(d-1)^\ell/(2\ell)$ is the limiting mean of the total
number of $\ell$-cycles in a random $d$-regular graph on $n$ vertices
as $n\to\infty$. It is a count for the whole graph, not a density, and
not an upper bound. The refutation of (i) uses $d=2$ only; for
$d\ge3$ claim (i) is true by Proposition~\ref{prop:upper}, so the
conjecture fails for $d\ge3$ through claim (ii) alone.
\end{remark}

\section*{The earlier submission}
Pull request 302 of the competition repository (by
\texttt{orionsheep}) treated this conjecture with the triangle $K_3$,
which is the case $\ell=3$ of Theorem~\ref{thm:cycle}. It was closed
without merging on 3 October 2026, together with a comment of its
author withdrawing the submission; its comment thread contains no
reviewer verdict. Its example and its conclusion were correct.

We read its Lean file. It modelled a finite graph by a structure of its
own and defined the cycle measure through maps from $\mathbb Z/n$ to
the vertices with cyclically adjacent consecutive values, identified by
their edge sets. These maps were not required to be injective. For
$n=3$ they are exactly the triangles, so the instance used was sound.
For $n\ge4$ they include closed walks that are not cycles, such as
$a,b,a,b$ on a single edge; with that definition a single edge would
have positive ``$4$-cycle measure''. So the universally quantified
statement negated there was not the bound for cycles in general. Its
statement about the tree was only the positivity of the expression
$(d-1)^\ell/(2\ell)$.

The present submission uses Mathlib's graphs and cycle graphs, counts
genuine cycles (injective homomorphisms from $C_\ell$), proves the
violation for every $\ell\ge3$, and separates what holds for $d\ge3$.

\section*{Formal verification and scope}
Graphs are Mathlib's \texttt{SimpleGraph}; degrees are Mathlib's
\texttt{SimpleGraph.degree}; the cycle graph is Mathlib's
\texttt{SimpleGraph.cycleGraph}. The type \texttt{CycleCopy G L}
consists of the injective graph homomorphisms from
\texttt{cycleGraph L} to \texttt{G}, and
\begin{center}\small
\texttt{cycleMeasure G L}
\end{center}
is their number divided by $2L$ and by the number of vertices, a
rational number. \texttt{claimedBound d L} is $(d-1)^L/(2L)$.
\texttt{rotation} and \texttt{reflection} are the maps of the proof of
Theorem~\ref{thm:cycle}, and \texttt{two\_mul\_le\_card} shows that
they are $2L$ different elements. The theorem
\begin{center}\small
\texttt{cycle\_violates}
\end{center}
states that \texttt{claimedBound 2 (l + 3)} is smaller than the cycle
measure of \texttt{cycleGraph (l + 3)} for every natural number $l$.

The proposition \texttt{ClaimedUpperBound} is claim (i) for all finite
graphs, all $d$ and all $L\ge3$. Its negation is
\begin{center}\small\texttt{conjecture\_false},\end{center}
and \texttt{conjecture\_false\_every\_length} gives, for each $L\ge3$,
a graph with all degrees at most $2$ violating the bound for $d=2$. No
\texttt{sorry}, \texttt{native\_decide} or custom axiom occurs.

Scope. Lean treats finite graphs, which are special graph limits;
general graph limits are not formalised, and a counterexample among
finite graphs refutes claim (i) for graph limits. Lean proves
$c_\ell(C_\ell)\ge1/\ell$; the equality, Proposition~\ref{prop:upper},
the discussion of claim (ii), and the remark are proved on paper only.
The refutation in Lean does not use the degenerate value $d=0$, for
which the expression $(d-1)^L/(2L)$ is negative when $L$ is odd.

\section*{Source and reproduction}
The exact bilingual source is included byte-for-byte as
\texttt{SOURCE.md}. Its repository path is
\begin{center}\small\texttt{conjectures/00000004413.md}.\end{center}
The project pins Lean 4.19.0 and Mathlib commit
\begin{center}\small\texttt{c44e0c8ee63ca166450922a373c7409c5d26b00b}.\end{center}
From \texttt{lean/}, run \texttt{lake build}, then
\begin{center}\small\texttt{lake env lean -DwarningAsError=true Main.lean}.\end{center}
The included public-Git manifest locks all dependency revisions. The
\texttt{verification/} directory records the fresh-build logs,
theorem-axiom reports, artifact hashes, review and final PDF
inspection. No supplementary numerical computation is required.
\end{document}
